\documentclass{amsart}
% For dealing with references we use the comment environment
\usepackage{verbatim}
\newenvironment{reference}{\comment}{\endcomment}
%\newenvironment{reference}{}{}
\newenvironment{slogan}{\comment}{\endcomment}
\newenvironment{history}{\comment}{\endcomment}

% For commutative diagrams we use Xy-pic
\usepackage[all]{xy}

% We use 2cell for 2-commutative diagrams.
\xyoption{2cell}
\UseAllTwocells

% We use multicol for the list of chapters between chapters
\usepackage{multicol}

% This is generally recommended for better output
\usepackage{lmodern}
\usepackage[T1]{fontenc}

% For cross-file-references

% Package for hypertext links:
\usepackage{hyperref}

% For any local file, say "hello.tex" you want to link to please

% Theorem environments.
%
\theoremstyle{plain}
\newtheorem{theorem}[subsection]{Theorem}
\newtheorem{proposition}[subsection]{Proposition}
\newtheorem{lemma}[subsection]{Lemma}

\theoremstyle{definition}
\newtheorem{definition}[subsection]{Definition}
\newtheorem{example}[subsection]{Example}
\newtheorem{exercise}[subsection]{Exercise}
\newtheorem{situation}[subsection]{Situation}

\theoremstyle{remark}
\newtheorem{remark}[subsection]{Remark}
\newtheorem{remarks}[subsection]{Remarks}

\numberwithin{equation}{subsection}

% Macros
%
\def\lim{\mathop{\mathrm{lim}}\nolimits}
\def\colim{\mathop{\mathrm{colim}}\nolimits}
\def\Spec{\mathop{\mathrm{Spec}}}
\def\Hom{\mathop{\mathrm{Hom}}\nolimits}
\def\Ext{\mathop{\mathrm{Ext}}\nolimits}
\def\SheafHom{\mathop{\mathcal{H}\!\mathit{om}}\nolimits}
\def\SheafExt{\mathop{\mathcal{E}\!\mathit{xt}}\nolimits}
\def\Sch{\mathit{Sch}}
\def\Mor{\mathop{\mathrm{Mor}}\nolimits}
\def\Ob{\mathop{\mathrm{Ob}}\nolimits}
\def\Sh{\mathop{\mathit{Sh}}\nolimits}
\def\NL{\mathop{N\!L}\nolimits}
\def\CH{\mathop{\mathrm{CH}}\nolimits}
\def\proetale{{pro\text{-}\acute{e}tale}}
\def\etale{{\acute{e}tale}}
\def\QCoh{\mathit{QCoh}}
\def\Ker{\mathop{\mathrm{Ker}}}
\def\Im{\mathop{\mathrm{Im}}}
\def\Coker{\mathop{\mathrm{Coker}}}
\def\Coim{\mathop{\mathrm{Coim}}}

% Boxtimes
%
\DeclareMathSymbol{\boxtimes}{\mathbin}{AMSa}{"02}

%
% Macros for moduli stacks/spaces
%
\def\QCohstack{\mathcal{QC}\!\mathit{oh}}
\def\Cohstack{\mathcal{C}\!\mathit{oh}}
\def\Spacesstack{\mathcal{S}\!\mathit{paces}}
\def\Quotfunctor{\mathrm{Quot}}
\def\Hilbfunctor{\mathrm{Hilb}}
\def\Curvesstack{\mathcal{C}\!\mathit{urves}}
\def\Polarizedstack{\mathcal{P}\!\mathit{olarized}}
\def\Complexesstack{\mathcal{C}\!\mathit{omplexes}}
% \Pic is the operator that assigns to X its picard group, usage \Pic(X)
% \Picardstack_{X/B} denotes the Picard stack of X over B
% \Picardfunctor_{X/B} denotes the Picard functor of X over B
\def\Pic{\mathop{\mathrm{Pic}}\nolimits}
\def\Picardstack{\mathcal{P}\!\mathit{ic}}
\def\Picardfunctor{\mathrm{Pic}}
\def\Deformationcategory{\mathcal{D}\!\mathit{ef}}

\setlength{\emergencystretch}{3em}
\begin{document}
\title[Stacks Project, Tag 0B73]{449 --- Gysin pullbacks for line-bundle sections commute on Chow homology}
\maketitle
\noindent Copyright (C) 2005--2025 Johan de Jong. Stacks Project authors. Source: \href{https://stacks.math.columbia.edu/tag/0B73}{Tag 0B73}.

\noindent Editorial difficulty note (not author text). Difficulty H2: The proof must compare intersection cycles when their supports meet the divisor centres improperly. Carefully normalized meromorphic sections reduce the difference to tame-symbol divisors, and one must show every resulting rational-equivalence support lies inside the common centre rather than merely vanishes after pushforward to the ambient scheme..

\noindent Editorial provenance note (not author text). History: excerpt from revision \texttt{a0444\allowbreak 6e57e\allowbreak c1fbc\allowbreak 252a8\allowbreak 71afc\allowbreak ec775\allowbreak 2fb28\allowbreak 07b14}, chow.tex, lines 5495--5570. Reference presentation was adapted; original-author context and core support blocks were retained or added. The mathematical wording is unchanged. Licensed under GNU FDL 1.2 or later; no invariant sections or cover texts. See ../licenses/COPYING. Context blocks reproduce author definitions and supporting results; they are not additional counted problems.

\begin{situation}
\label{situation-setup}
Here $S$ is a locally Noetherian, and universally catenary scheme.
Moreover, we assume $S$ is endowed with a dimension function
$\delta : S \longrightarrow \mathbf{Z}$.
\end{situation}

\begin{definition}
\label{definition-gysin-homomorphism}
Let $(S, \delta)$ be as in Situation \ref{situation-setup}.
Let $X$ be locally of finite type over $S$.
Let $(\mathcal{L}, s)$ be a pair consisting of an invertible
sheaf and a global section $s \in \Gamma(X, \mathcal{L})$.
Let $D = Z(s)$ be the zero scheme of $s$, and
denote $i : D \to X$ the closed immersion.
We define, for every integer $k$, a {\it Gysin homomorphism}
$$
i^* : Z_{k + 1}(X) \to \CH_k(D).
$$
by the following rules:
\begin{enumerate}
\item Given an integral closed subscheme $W \subset X$ with
$\dim_\delta(W) = k + 1$ we define
\begin{enumerate}
\item if $W \not \subset D$, then $i^*[W] = [D \cap W]_k$ as a
$k$-cycle on $D$, and
\item if $W \subset D$, then
$i^*[W] = i'_*(c_1(\mathcal{L}|_W) \cap [W])$,
where $i' : W \to D$ is the induced closed immersion.
\end{enumerate}
\item For a general $(k + 1)$-cycle $\alpha = \sum n_j[W_j]$
we set
$$
i^*\alpha = \sum n_j i^*[W_j]
$$
\item If $D$ is an effective Cartier divisor, then we denote
$D \cdot \alpha = i_*i^*\alpha$ the pushforward of the class $i^*\alpha$
to a class on $X$.
\end{enumerate}
\end{definition}

\begin{lemma}
\label{lemma-gysin-commutes-gysin}
Let $(S, \delta)$ be as in Situation \ref{situation-setup}. Let $X$ be locally
of finite type over $S$. Let $(\mathcal{L}, s, i : D \to X)$ and
$(\mathcal{L}', s', i' : D' \to X)$ be two triples as in
Definition \ref{definition-gysin-homomorphism}. Then the diagram
$$
\xymatrix{
\CH_k(X) \ar[r]_{i^*} \ar[d]_{(i')^*} & \CH_{k - 1}(D) \ar[d]^{j^*} \\
\CH_{k - 1}(D') \ar[r]^{(j')^*} & \CH_{k - 2}(D \cap D')
}
$$
commutes where each of the maps is a Gysin map.
\end{lemma}

\begin{proof}
Denote $j : D \cap D' \to D$ and $j' : D \cap D' \to D'$ the closed
immersions corresponding to $(\mathcal{L}|_{D'}, s|_{D'}$ and
$(\mathcal{L}'_D, s|_D)$. We have to show that
$(j')^*i^*\alpha = j^* (i')^*\alpha$ for all $\alpha \in \CH_k(X)$.
Let $W \subset X$ be an integral closed subscheme of dimension $k$.
Let us prove the equality in case $\alpha = [W]$. We will deduce
it from the key formula.

\medskip\noindent
We let $\sigma$ be a nonzero meromorphic section of $\mathcal{L}|_W$
which we require to be equal to $s|_W$ if $W \not \subset D$.
We let $\sigma'$ be a nonzero meromorphic section of $\mathcal{L}'|_W$
which we require to be equal to $s'|_W$ if $W \not \subset D'$.
Write
$$
\text{div}_{\mathcal{L}|_W}(\sigma) =
\sum \text{ord}_{Z_i, \mathcal{L}|_W}(\sigma)[Z_i] = \sum n_i[Z_i]
$$
and similarly
$$
\text{div}_{\mathcal{L}'|_W}(\sigma') =
\sum \text{ord}_{Z_i, \mathcal{L}'|_W}(\sigma')[Z_i] = \sum n'_i[Z_i]
$$
as in the discussion in Section \href{https://stacks.math.columbia.edu/tag/0AYB}{section key (Tag 0AYB)}.
Then we see that $Z_i \subset D$ if $n_i \not = 0$ and
$Z'_i \subset D'$ if $n'_i \not = 0$. For each $i$, let $\xi_i \in Z_i$
be the generic point. As in Section \href{https://stacks.math.columbia.edu/tag/0AYB}{section key (Tag 0AYB)} we choose
for each $i$ an element
$\sigma_i \in \mathcal{L}_{\xi_i}$, resp.\ $\sigma'_i \in \mathcal{L}'_{\xi_i}$
which generates over $B_i = \mathcal{O}_{W, \xi_i}$
and which is equal to the image of
$s$, resp.\ $s'$ if $Z_i \not \subset D$, resp.\ $Z_i \not \subset D'$.
Write $\sigma = f_i \sigma_i$ and $\sigma' = f'_i\sigma'_i$ so that
$n_i = \text{ord}_{B_i}(f_i)$ and
$n'_i = \text{ord}_{B_i}(f'_i)$.
From our definitions it follows that
$$
(j')^*i^*[W] =
\sum \text{ord}_{B_i}(f_i) \text{div}_{\mathcal{L}'|_{Z_i}}(\sigma'_i|_{Z_i})
$$
as cycles and
$$
j^*(i')^*[W] =
\sum \text{ord}_{B_i}(f'_i) \text{div}_{\mathcal{L}|_{Z_i}}(\sigma_i|_{Z_i})
$$
The key formula (Lemma \href{https://stacks.math.columbia.edu/tag/0AYC}{lemma key formula (Tag 0AYC)}) now gives the equality
$$
\sum \left(
\text{ord}_{B_i}(f_i) \text{div}_{\mathcal{L}'|_{Z_i}}(\sigma'_i|_{Z_i}) -
\text{ord}_{B_i}(f'_i) \text{div}_{\mathcal{L}|_{Z_i}}(\sigma_i|_{Z_i})
\right) =
\sum \text{div}_{Z_i}(\partial_{B_i}(f_i, f'_i))
$$
of cycles. Note that $\text{div}_{Z_i}(\partial_{B_i}(f_i, f'_i)) = 0$ if
$Z_i \not \subset D \cap D'$ because in this case either $f_i = 1$
or $f'_i = 1$. Thus we get a rational equivalence between our specific
cycles representing $(j')^*i^*[W]$ and $j^*(i')^*[W]$ on $D \cap D' \cap W$.
By Remark \href{https://stacks.math.columbia.edu/tag/02RZ}{remark infinite sums rational equivalences (Tag 02RZ)}
the result follows for general $\alpha$.
\end{proof}
\end{document}
